% TOP_PROBLEM: {"id": 3329, "title": "Growth of finitely presented groups", "category_id": 17, "category": {"id": 17, "name": "group_theory", "display_name": "Group Theory", "description": "Problems about groups, group actions, representations, and related algebraic structures.", "slug": "group-theory", "order_index": 17, "created_at": "2026-07-31T15:26:25.671Z"}, "status": "open"}
% FIRST_POSTED: 2026-09-27
% !TeX program = lualatex
% ATTEMPT_STATUS: unresolved
\documentclass[11pt]{article}
\usepackage[margin=25mm]{geometry}
\usepackage{fontspec}
\setmainfont{Latin Modern Roman}
\usepackage{amsmath,amssymb,mathtools}
\usepackage{unicode-math}
\setmathfont{Latin Modern Math}
\usepackage{xurl,hyperref}
\hypersetup{colorlinks=true,urlcolor=blue}
\setcounter{secnumdepth}{0}
\setlength{\parindent}{0pt}
\setlength{\parskip}{5pt}
\setlength{\emergencystretch}{3em}
\begin{document}
\section*{\texorpdfstring{Growth of finitely presented groups}{Growth of finitely presented groups}}\label{3329}
\subsection{Definitions and mathematical statement}
For a finitely generated group $G$ with finite symmetric generating set $S$, let
\[
 \gamma_{G,S}(n)=\#\{g\in G:|g|_S\le n\},
\]
where $|g|_S$ is the minimum word length. Intermediate growth means simultaneously
\[
 \lim_{n\to\infty}\frac{\log\gamma_{G,S}(n)}n=0,
 \qquad\text{and}\qquad
 \nexists C,d:\gamma_{G,S}(n)\le C(n+1)^d\text{ for all }n.
\]
The existence problem asks for such a $G$ that also has a finite presentation, with finitely many defining relations. These growth properties are unchanged on changing the finite generating set. No precise intermediate rate is prescribed, and an infinitely related finitely generated construction does not meet the extra finite-presentation requirement.
\par\smallskip\textit{Formulation and concept references:} \href{https://ar5iv.labs.arxiv.org/html/1111.0512}{[S1]}.

\subsection{Short English statement}
Does a finitely presented group exist whose growth is slower than exponential but faster than every polynomial bound?
\subsection{Research attempt}
\paragraph{Scope and source audit (27 September 2026).}
The requested example must satisfy three separate conditions: finite
generation, finite presentation, and the two growth conditions in the
statement. A finite recursive description or a finite-state automaton
is not the same as a finite list of defining relations. The source
survey [S1] explicitly distinguishes these notions and the known
intermediate-growth constructions from the finite-presentation problem.
Its full text was inspected; no general resolution is imported here.

The phrase ``faster than every polynomial'' is used in the precise
sense printed in the statement: no polynomial is a global upper bound.
We do not silently replace it by a stronger regularity assertion about
a limit for every quotient $\gamma(n)/n^d$. Growth functions may have
substantial changes of scale, so all estimates below keep their
quantifiers visible.

The attack examines two possible routes: truncate the relations of a
known intermediate-growth group, or place it inside a finitely presented
amenable group. Exact growth comparisons and explicit examples show
why neither route gives the required conclusion without additional
uniform estimates. We also prove the amenability consequence of
subexponential growth and a useful local-stabilization theorem for
increasing relator lists.

\paragraph{Submultiplicativity gives an actual exponential growth rate.}
Let $B(n)$ be the word ball and $\gamma(n)=|B(n)|$. Every word of
length at most $n+m$ splits into a prefix of length at most $n$ and
a suffix of length at most $m$, so
\[
 B(n+m)\subseteq B(n)B(m),\qquad
 \gamma(n+m)\le\gamma(n)\gamma(m).
 \tag{374.1}
\]
Put $a_n=\log\gamma(n)$. Then $a_{n+m}\le a_n+a_m$.
A direct proof of the subadditive limit formula is short. For a fixed
$k\ge1$, write $n=qk+r$, $0\le r<k$. Then
$a_n\le q a_k+a_r$, so
$\limsup a_n/n\le a_k/k$. Taking the infimum over $k$, and comparing
with the automatic lower bound by that infimum, gives
\[
 h(G,S)=\lim_{n\to\infty}\frac{\log\gamma(n)}n
       =\inf_{n\ge1}\frac{\log\gamma(n)}n\ge0.
 \tag{374.2}
\]
Thus the first growth condition is equivalent to $h(G,S)=0$.
For each $\varepsilon>0$, it implies
$\gamma(n)\le C_\varepsilon e^{\varepsilon n}$ for all $n$, with
a finite constant absorbing the initial segment. It does not give a
single polynomial bound: the constant and the cutoff may deteriorate
as $\varepsilon\downarrow0$.

\paragraph{Changing generators preserves both requested conditions.}
For two finite symmetric generating sets $S,T$, there are finite
constants $L,L'$ such that each generator of $S$ has $T$-length at
most $L$ and each generator of $T$ has $S$-length at most $L'$.
Replacing letters in words gives
\[
 \gamma_{G,S}(n)\le\gamma_{G,T}(Ln),\qquad
 \gamma_{G,T}(n)\le\gamma_{G,S}(L'n).
 \tag{374.3}
\]
If one growth rate is zero, these inequalities force the other to be
zero by (374.2). A polynomial upper bound on either side similarly
transfers to the other after changing its constant. Hence absence of
such a bound also transfers. No equivalence of exact pointwise formulas
or of leading constants is being asserted.

The same argument gives useful one-sided comparisons for maps.
A quotient generated by the images of $S$ has at most $\gamma_{G,S}(n)$
elements in its radius-$n$ ball. A finitely generated subgroup with
generators of ambient length at most $L$ has intrinsic growth at most
$\gamma_{G,S}(Ln)$. The latter says that a group containing an
intermediate-growth subgroup cannot have polynomial growth. It gives
no subexponential \emph{upper} bound on the larger group.

\paragraph{Finite-index passage can be controlled on both sides.}
Let $H\le G$ have finite index $m$, and choose representatives $r$
for the right cosets $Hg$, including $1$. Let each representative have
$S$-length at most $L$. The finite Schreier generating set for $H$
consists of the nonidentity elements
\[
 t=r s(r')^{-1},\qquad Hr s=Hr',\quad s\in S.
\]
It can be taken symmetric and each such generator has ambient length
at most $2L+1$. Reading a word one letter at a time rewrites its value
as a product of at most its length many Schreier generators followed
by one representative. For elements of $H$ the final representative
is $1$. Thus these elements really generate $H$, and
\[
 \gamma_{H,T}(n)\le\gamma_{G,S}((2L+1)n),\qquad
 \gamma_{G,S}(n)\le m\gamma_{H,T}(n).
 \tag{374.4}
\]
The second inequality follows by counting the final representative and
the preceding Schreier word; uniqueness is not needed for this upper
bound. Subexponential growth and polynomial boundedness therefore
pass in both directions across finite index. This gives a valid way
to move an example between commensurable groups, but does not turn an
arbitrary infinite-index embedding into such a comparison.

\paragraph{Two fully explicit finitely presented boundary classes.}
The group $\mathbb Z^d$ has the finite presentation with $d$ generators
and all pairwise commutator relations. With the standard generators,
word length is the $\ell^1$ norm. Hence for $d\ge1$,
\[
 (2\lfloor n/d\rfloor+1)^d
       \le\gamma_{\mathbb Z^d}(n)\le(2n+1)^d.
\]
This is polynomial growth, so it fails the second requested condition.
Finite groups also fail it, since their growth is bounded.

The free group of rank $r\ge2$ has a finite presentation with no
relators. Its nonidentity reduced words of length $j$ number
$2r(2r-1)^{j-1}$: choose the first letter freely and forbid only the
inverse of the preceding letter thereafter. Thus
\[
 \gamma_{F_r}(n)=1+2r\sum_{j=1}^n(2r-1)^{j-1},
 \qquad h(F_r,S)=\log(2r-1)>0.
\]
It fails the first condition. These elementary examples show that a
finite relator list by itself allows both familiar extremes, while
locating neither an intermediate construction nor an obstruction to it.

\paragraph{Subexponential growth supplies Følner sets.}
\textbf{Proposition 374.1.} Every finitely generated group with
$h(G,S)=0$ is amenable.

\emph{Proof.} For every $\eta>0$ there are arbitrarily large integers
$n$ with
$\gamma(n+1)\le(1+\eta)\gamma(n)$. Otherwise for all sufficiently
large $n$ the ratio would exceed $1+\eta$, and iterating would give
positive exponential rate, contrary to (374.2).

For a generator $s$, both $sB(n)$ and $B(n)$ lie in $B(n+1)$ and
have equal cardinality. Therefore
\[
 \frac{|sB(n)\mathbin\triangle B(n)|}{|B(n)|}
 \le2\frac{\gamma(n+1)-\gamma(n)}{\gamma(n)}.
 \tag{374.5}
\]
Choose $n_j\to\infty$ with this ratio of consecutive ball sizes
at most $1+1/j$. The normalized averages on $B(n_j)$ are positive
normalized functionals on $\ell^\infty(G)$. By weak-star compactness
of the dual unit ball, take a convergent subnet. Estimate (374.5)
makes its limit invariant under every generator, hence under every
finite product of generators. The limit is an invariant mean.
\hfill$\square$

Only a sequence of good radii is needed. We did not assume that
$\gamma(n+1)/\gamma(n)$ tends to one at every radius. This is an
important distinction for attempts based on observed local growth.
The weak-star compactness passage proves existence of a mean, not an
efficient algorithm for testing amenability from a presentation.

\paragraph{Amenability alone leaves exponential growth possible.}
Consider the lamplighter group
$L=(\bigoplus_{\mathbb Z}C_2)\rtimes\mathbb Z$. Write elements as
$(f,k)$, with finitely supported binary functions $f$, and multiply by
\[
 (f,k)(g,l)=(f+\operatorname{shift}_k g,k+l).
\]
Let $a$ toggle the lamp at zero and let $t$ move the cursor one step.
For each binary string $(\varepsilon_0,\ldots,\varepsilon_{m-1})$,
the word
\[
 a^{\varepsilon_0}t\,a^{\varepsilon_1}t\cdots
                         a^{\varepsilon_{m-1}}t
\]
has length at most $2m$, ends with cursor $m$, and records those
$m$ lamp choices at distinct sites. The resulting $2^m$ elements
are distinct. Thus $\gamma_L(2m)\ge2^m$, implying positive
exponential growth rate by monotonicity and (374.2).

Nevertheless it is amenable by explicit Følner sets. Let $F_m$ consist
of configurations supported in $[-m,m]$ with cursor also in $[-m,m]$.
Its size is $(2m+1)2^{2m+1}$. Right multiplication by $a$ preserves
$F_m$, because it toggles the lamp at the current cursor. Right
multiplication by $t$ or $t^{-1}$ changes only the allowed cursor
interval, so
\[
 |F_m t\mathbin\triangle F_m|/|F_m|=2/(2m+1).
\]
Right averages therefore give a right-invariant mean, and inversion
converts it to a left-invariant one. This proves amenability directly.

The example is used only to disprove the implication ``amenable implies
subexponential''. It is not offered as a finitely presented candidate.
In particular a finite-state or finite-generator description of this
wreath product is not treated as a finite presentation.

\paragraph{A finite presentation can already permit an exponential semigroup.}
For another exact growth test, take
\[
 B=\langle a,t\mid tat^{-1}=a^2\rangle.
\]

Define affine transformations of the real line by $a(x)=x+1$ and
$t(x)=2x$. They satisfy the relation, so give a homomorphism from $B$
to the affine group. The $2^m$ words
\[
 a^{\varepsilon_0}t\,a^{\varepsilon_1}t\cdots
                         a^{\varepsilon_{m-1}}t
\]
act as $x\mapsto2^m x+b$, where
$b=\sum_{i=0}^{m-1}2^i\varepsilon_i$ under right-to-left composition.
They have distinct translation coefficients and length at most $2m$.
Thus this finitely presented group also has exponential growth. Faithfulness
of the entire affine representation is not needed: different images
already certify different group elements.

This calculation is a useful adversarial test for constructions involving
an ascending conjugation relation. Such a relation can encode a finite
recursive scheme, but it may also create exponentially many independently
addressable choices. Compressing infinitely many descriptions into a
short presentation is not automatically compatible with a subexponential
upper bound.

\paragraph{Truncating an infinite relation list reverses the useful upper bound.}
Suppose a finitely generated group is given by an infinite presentation
\[
 G=F(S)/\langle\!\langle R_1,R_2,\ldots\rangle\!\rangle,
 \qquad G_m=F(S)/\langle\!\langle R_1,\ldots,R_m\rangle\!\rangle.
\]
Every $G_m$ is finitely presented and maps onto $G$. Consequently
\[
 \gamma_G(n)\le\gamma_{G_m}(n)
 \quad\text{for every }m,n.
 \tag{374.6}
\]
If $G$ has no polynomial growth bound, neither does any $G_m$; otherwise
that bound would descend to $G$. This successfully transfers the
superpolynomial side of the desired example. It transfers the
subexponential side in the wrong direction: an upper bound for $G$
says nothing about an upper bound for its larger approximants $G_m$.

There is a genuine local stabilization theorem, but its quantifiers
are weaker than the needed global estimate.
\textbf{Lemma 374.2.} For each fixed radius $R$, there exists $m_R$
such that the marked radius-$R$ balls in $G_m$ and $G$ agree for all
$m\ge m_R$.

\emph{Proof.} Two words of length at most $R+1$ are equal in $G$
exactly when their quotient word, of length at most $2R+2$, belongs
to the normal closure of the whole relator set. Each such membership
has a finite expression as a product of conjugates of relators and
inverses. Only finitely many relators occur in that expression.
There are finitely many pairs of words of length at most $R+1$.
Choose $m_R$ large enough to include the relators in an expression
for every pair that is equal in $G$. Every such equality then holds
in $G_m$. Conversely every equality in $G_m$ holds in its quotient
$G$. This proves equality of the ball identifications and their marked
edge incidences wherever the edges remain in the ball: an edge from
$u$ to $v$ labeled $s$ tests the equality $us=v$, and $us$ has
length at most $R+1$.\hfill$\square$

The lemma does not give one finite $m$ that works at all radii. Even
if $m_R$ were computable, replacing $m$ by a radius-dependent value
would produce a sequence of groups rather than one finitely presented
group with the required growth estimate.

\paragraph{Exactly which uniform estimate would make truncation succeed.}
Suppose $G$ has intermediate growth and, for some fixed $m$, one can
prove
\[
 \forall\varepsilon>0\ \exists C_\varepsilon<\infty\ 
 \forall n:\qquad\gamma_{G_m}(n)\le C_\varepsilon e^{\varepsilon n}.
 \tag{374.7}
\]
Then (374.2) gives subexponential growth of $G_m$, while (374.6)
excludes any polynomial upper bound for it. Since $G_m$ is finitely
presented, this would solve the existence problem. The entire difficulty
of that route is therefore concentrated in proving (374.7) for one
fixed truncation, not in the formal passage to a finite presentation.

A family of estimates only for $n\le R_m$ with $R_m\to\infty$ does
not have the quantifier order in (374.7). Nor does an estimate with
exponential rate $\varepsilon_m\downarrow0$ give a member of rate
zero: every $\varepsilon_m$ may remain positive. A limit group can
have zero rate even though each fixed approximant still grows
exponentially. No continuity theorem ruling out this possibility is
proved or assumed here.

\paragraph{Subgroup embeddings and quotients have different failures.}
If an intermediate-growth group embeds into a finitely presented group
$H$, the comparison after (374.3) rules out polynomial growth of $H$.
But $H$ can have many more elements in its balls, so subexponential
upper bounds do not transfer upward. Assuming $H$ amenable only
provides the insufficient condition exposed by the lamplighter example.
Additional distortion estimates would control the embedded subgroup's
metric but would still not count all elements outside it.

In the other direction, taking a quotient of a finitely presented group
preserves finite presentation only when the new normal subgroup is
finitely normally generated. An infinitely related quotient can have
the desired growth without satisfying this finiteness condition. The
normal-closure expression for each individual word is finite, but a
single finite list generating \emph{all} extra relations is a different
assertion. Lemma 374.2 makes the distinction explicit.

\paragraph{Diagnostics and outcome.}
Exact finite checks count reduced free words, integer lattice balls,
and the distinct affine images in the exponential diagnostic. They
also enumerate lamp configurations to verify the stated Følner boundary
ratios. These tests establish finite formulas and expose specific failed
implications; no finite collection of growth values determines whether
one presentation satisfies both asymptotic conditions.

\textbf{Outcome: UNRESOLVED.} The attempt proves the generator and
finite-index invariances, the amenability consequence, and the local
stabilization lemma with its exact quantifiers. It identifies a sufficient
uniform bound for a finite truncation and why embeddings or amenability
do not replace that bound. No fixed finitely presented group with a
proved zero exponential rate and absence of every polynomial upper
bound is constructed, and no general impossibility theorem is obtained.

\subsection{Sources}
\begin{itemize}
\item[S1]R. Grigorchuk, Milnor's Problem on the Growth of Groups and its Consequences, arXiv:1111.0512 (2011, revised 2013).
\url{https://ar5iv.labs.arxiv.org/html/1111.0512}.
\textit{Formulation source.}
\end{itemize}
\end{document}
